\IfFileExists{stacks-project.cls}{%
\documentclass{stacks-project}
}{%
\documentclass{amsart}
}

% For dealing with references we use the comment environment
\usepackage{verbatim}
\newenvironment{reference}{\comment}{\endcomment}
%\newenvironment{reference}{}{}
\newenvironment{slogan}{\comment}{\endcomment}
\newenvironment{history}{\comment}{\endcomment}

% For commutative diagrams we use Xy-pic
\usepackage[all]{xy}

% We use 2cell for 2-commutative diagrams.
\xyoption{2cell}
\UseAllTwocells

% We use multicol for the list of chapters between chapters
\usepackage{multicol}

% This is generally recommended for better output
\usepackage{lmodern}
\usepackage[T1]{fontenc}

% For cross-file-references
\usepackage{xr-hyper}

% Package for hypertext links:
\usepackage{hyperref}

% For any local file, say "hello.tex" you want to link to please
% use \externaldocument[hello-]{hello}
\externaldocument[introduction-]{introduction}
\externaldocument[conventions-]{conventions}
\externaldocument[sets-]{sets}
\externaldocument[categories-]{categories}
\externaldocument[topology-]{topology}
\externaldocument[sheaves-]{sheaves}
\externaldocument[sites-]{sites}
\externaldocument[stacks-]{stacks}
\externaldocument[fields-]{fields}
\externaldocument[algebra-]{algebra}
\externaldocument[brauer-]{brauer}
\externaldocument[homology-]{homology}
\externaldocument[derived-]{derived}
\externaldocument[simplicial-]{simplicial}
\externaldocument[more-algebra-]{more-algebra}
\externaldocument[smoothing-]{smoothing}
\externaldocument[modules-]{modules}
\externaldocument[sites-modules-]{sites-modules}
\externaldocument[injectives-]{injectives}
\externaldocument[cohomology-]{cohomology}
\externaldocument[sites-cohomology-]{sites-cohomology}
\externaldocument[dga-]{dga}
\externaldocument[dpa-]{dpa}
\externaldocument[sdga-]{sdga}
\externaldocument[hypercovering-]{hypercovering}
\externaldocument[schemes-]{schemes}
\externaldocument[constructions-]{constructions}
\externaldocument[properties-]{properties}
\externaldocument[morphisms-]{morphisms}
\externaldocument[coherent-]{coherent}
\externaldocument[divisors-]{divisors}
\externaldocument[limits-]{limits}
\externaldocument[varieties-]{varieties}
\externaldocument[topologies-]{topologies}
\externaldocument[descent-]{descent}
\externaldocument[perfect-]{perfect}
\externaldocument[more-morphisms-]{more-morphisms}
\externaldocument[flat-]{flat}
\externaldocument[groupoids-]{groupoids}
\externaldocument[more-groupoids-]{more-groupoids}
\externaldocument[etale-]{etale}
\externaldocument[chow-]{chow}
\externaldocument[intersection-]{intersection}
\externaldocument[pic-]{pic}
\externaldocument[weil-]{weil}
\externaldocument[adequate-]{adequate}
\externaldocument[dualizing-]{dualizing}
\externaldocument[duality-]{duality}
\externaldocument[discriminant-]{discriminant}
\externaldocument[derham-]{derham}
\externaldocument[local-cohomology-]{local-cohomology}
\externaldocument[algebraization-]{algebraization}
\externaldocument[curves-]{curves}
\externaldocument[resolve-]{resolve}
\externaldocument[models-]{models}
\externaldocument[functors-]{functors}
\externaldocument[equiv-]{equiv}
\externaldocument[pione-]{pione}
\externaldocument[etale-cohomology-]{etale-cohomology}
\externaldocument[proetale-]{proetale}
\externaldocument[relative-cycles-]{relative-cycles}
\externaldocument[more-etale-]{more-etale}
\externaldocument[trace-]{trace}
\externaldocument[crystalline-]{crystalline}
\externaldocument[spaces-]{spaces}
\externaldocument[spaces-properties-]{spaces-properties}
\externaldocument[spaces-morphisms-]{spaces-morphisms}
\externaldocument[decent-spaces-]{decent-spaces}
\externaldocument[spaces-cohomology-]{spaces-cohomology}
\externaldocument[spaces-limits-]{spaces-limits}
\externaldocument[spaces-divisors-]{spaces-divisors}
\externaldocument[spaces-over-fields-]{spaces-over-fields}
\externaldocument[spaces-topologies-]{spaces-topologies}
\externaldocument[spaces-descent-]{spaces-descent}
\externaldocument[spaces-perfect-]{spaces-perfect}
\externaldocument[spaces-more-morphisms-]{spaces-more-morphisms}
\externaldocument[spaces-flat-]{spaces-flat}
\externaldocument[spaces-groupoids-]{spaces-groupoids}
\externaldocument[spaces-more-groupoids-]{spaces-more-groupoids}
\externaldocument[bootstrap-]{bootstrap}
\externaldocument[spaces-pushouts-]{spaces-pushouts}
\externaldocument[spaces-chow-]{spaces-chow}
\externaldocument[groupoids-quotients-]{groupoids-quotients}
\externaldocument[spaces-more-cohomology-]{spaces-more-cohomology}
\externaldocument[spaces-simplicial-]{spaces-simplicial}
\externaldocument[spaces-duality-]{spaces-duality}
\externaldocument[formal-spaces-]{formal-spaces}
\externaldocument[restricted-]{restricted}
\externaldocument[spaces-resolve-]{spaces-resolve}
\externaldocument[formal-defos-]{formal-defos}
\externaldocument[defos-]{defos}
\externaldocument[cotangent-]{cotangent}
\externaldocument[examples-defos-]{examples-defos}
\externaldocument[algebraic-]{algebraic}
\externaldocument[examples-stacks-]{examples-stacks}
\externaldocument[stacks-sheaves-]{stacks-sheaves}
\externaldocument[criteria-]{criteria}
\externaldocument[artin-]{artin}
\externaldocument[quot-]{quot}
\externaldocument[stacks-properties-]{stacks-properties}
\externaldocument[stacks-morphisms-]{stacks-morphisms}
\externaldocument[stacks-limits-]{stacks-limits}
\externaldocument[stacks-cohomology-]{stacks-cohomology}
\externaldocument[stacks-perfect-]{stacks-perfect}
\externaldocument[stacks-introduction-]{stacks-introduction}
\externaldocument[stacks-more-morphisms-]{stacks-more-morphisms}
\externaldocument[stacks-geometry-]{stacks-geometry}
\externaldocument[moduli-]{moduli}
\externaldocument[moduli-curves-]{moduli-curves}
\externaldocument[examples-]{examples}
\externaldocument[exercises-]{exercises}
\externaldocument[guide-]{guide}
\externaldocument[desirables-]{desirables}
\externaldocument[coding-]{coding}
\externaldocument[obsolete-]{obsolete}
\externaldocument[fdl-]{fdl}
\externaldocument[index-]{index}

% Theorem environments.
%
\theoremstyle{plain}
\newtheorem{theorem}[subsection]{Theorem}
\newtheorem{proposition}[subsection]{Proposition}
\newtheorem{lemma}[subsection]{Lemma}

\theoremstyle{definition}
\newtheorem{definition}[subsection]{Definition}
\newtheorem{example}[subsection]{Example}
\newtheorem{exercise}[subsection]{Exercise}
\newtheorem{situation}[subsection]{Situation}

\theoremstyle{remark}
\newtheorem{remark}[subsection]{Remark}
\newtheorem{remarks}[subsection]{Remarks}

\numberwithin{equation}{subsection}

% Macros
%
\def\lim{\mathop{\mathrm{lim}}\nolimits}
\def\colim{\mathop{\mathrm{colim}}\nolimits}
\def\Spec{\mathop{\mathrm{Spec}}}
\def\Hom{\mathop{\mathrm{Hom}}\nolimits}
\def\Ext{\mathop{\mathrm{Ext}}\nolimits}
\def\SheafHom{\mathop{\mathcal{H}\!\mathit{om}}\nolimits}
\def\SheafExt{\mathop{\mathcal{E}\!\mathit{xt}}\nolimits}
\def\Sch{\mathit{Sch}}
\def\Mor{\mathop{\mathrm{Mor}}\nolimits}
\def\Ob{\mathop{\mathrm{Ob}}\nolimits}
\def\Sh{\mathop{\mathit{Sh}}\nolimits}
\def\NL{\mathop{N\!L}\nolimits}
\def\CH{\mathop{\mathrm{CH}}\nolimits}
\def\proetale{{pro\text{-}\acute{e}tale}}
\def\etale{{\acute{e}tale}}
\def\QCoh{\mathit{QCoh}}
\def\Ker{\mathop{\mathrm{Ker}}}
\def\Im{\mathop{\mathrm{Im}}}
\def\Coker{\mathop{\mathrm{Coker}}}
\def\Coim{\mathop{\mathrm{Coim}}}

% Boxtimes
%
\DeclareMathSymbol{\boxtimes}{\mathbin}{AMSa}{"02}

%
% Macros for moduli stacks/spaces
%
\def\QCohstack{\mathcal{QC}\!\mathit{oh}}
\def\Cohstack{\mathcal{C}\!\mathit{oh}}
\def\Spacesstack{\mathcal{S}\!\mathit{paces}}
\def\Quotfunctor{\mathrm{Quot}}
\def\Hilbfunctor{\mathrm{Hilb}}
\def\Curvesstack{\mathcal{C}\!\mathit{urves}}
\def\Polarizedstack{\mathcal{P}\!\mathit{olarized}}
\def\Complexesstack{\mathcal{C}\!\mathit{omplexes}}
% \Pic is the operator that assigns to X its picard group, usage \Pic(X)
% \Picardstack_{X/B} denotes the Picard stack of X over B
% \Picardfunctor_{X/B} denotes the Picard functor of X over B
\def\Pic{\mathop{\mathrm{Pic}}\nolimits}
\def\Picardstack{\mathcal{P}\!\mathit{ic}}
\def\Picardfunctor{\mathrm{Pic}}
\def\Deformationcategory{\mathcal{D}\!\mathit{ef}}
\usepackage{cleveref}
\crefname{originalreference}{original reference}{original references}
\Crefname{originalreference}{Original reference}{Original references}
\makeatletter
\newlabel{algebra-lemma-flat-increases-depth}{{00LM}{source}{}{sourcepointer-1}{}}
\newlabel{algebra-lemma-flat-increases-depth@cref}{{[originalreference][0][]00LM}{[1][0][]source}{}{}{}}
\makeatother
\crefname{originalreference}{original reference}{original references}
\Crefname{originalreference}{Original reference}{Original references}
\makeatletter
\newlabel{algebra-lemma-grothendieck-regular-sequence}{{00MG}{source}{}{sourcepointer-2}{}}
\newlabel{algebra-lemma-grothendieck-regular-sequence@cref}{{[originalreference][0][]00MG}{[1][0][]source}{}{}{}}
\makeatother
\crefname{originalreference}{original reference}{original references}
\Crefname{originalreference}{Original reference}{Original references}
\makeatletter
\newlabel{algebra-lemma-regular-sequence-in-neighbourhood}{{061L}{source}{}{sourcepointer-3}{}}
\newlabel{algebra-lemma-regular-sequence-in-neighbourhood@cref}{{[originalreference][0][]061L}{[1][0][]source}{}{}{}}
\makeatother
\crefname{originallemma}{original lemma}{original lemmas}
\Crefname{originallemma}{Original lemma}{Original lemmas}
\makeatletter
\newlabel{lemma-base-change-rig-smooth-continuous}{{0GCL}{source}{}{sourcepointer-4}{}}
\newlabel{lemma-base-change-rig-smooth-continuous@cref}{{[originallemma][0][]0GCL}{[1][0][]source}{}{}{}}
\makeatother
\crefname{originallemma}{original lemma}{original lemmas}
\Crefname{originallemma}{Original lemma}{Original lemmas}
\makeatletter
\newlabel{lemma-equivalent-with-artin-smooth}{{0GAJ}{source}{}{sourcepointer-5}{}}
\newlabel{lemma-equivalent-with-artin-smooth@cref}{{[originallemma][0][]0GAJ}{[1][0][]source}{}{}{}}
\makeatother
\crefname{originallemma}{original lemma}{original lemmas}
\Crefname{originallemma}{Original lemma}{Original lemmas}
\makeatletter
\newlabel{lemma-fibre-regular}{{0GGJ}{source}{}{sourcepointer-6}{}}
\newlabel{lemma-fibre-regular@cref}{{[originallemma][0][]0GGJ}{[1][0][]source}{}{}{}}
\makeatother
\crefname{originallemma}{original lemma}{original lemmas}
\Crefname{originallemma}{Original lemma}{Original lemmas}
\makeatletter
\newlabel{lemma-rig-closed-point-after-localization}{{0GGG}{source}{}{sourcepointer-7}{}}
\newlabel{lemma-rig-closed-point-after-localization@cref}{{[originallemma][0][]0GGG}{[1][0][]source}{}{}{}}
\makeatother
\crefname{originallemma}{original lemma}{original lemmas}
\Crefname{originallemma}{Original lemma}{Original lemmas}
\makeatletter
\newlabel{lemma-rig-smooth-axioms}{{0GCK}{source}{}{sourcepointer-8}{}}
\newlabel{lemma-rig-smooth-axioms@cref}{{[originallemma][0][]0GCK}{[1][0][]source}{}{}{}}
\makeatother
\crefname{originalreference}{original reference}{original references}
\Crefname{originalreference}{Original reference}{Original references}
\makeatletter
\newlabel{more-algebra-lemma-flat-descent-regular-ideal}{{068N}{source}{}{sourcepointer-9}{}}
\newlabel{more-algebra-lemma-flat-descent-regular-ideal@cref}{{[originalreference][0][]068N}{[1][0][]source}{}{}{}}
\makeatother
\crefname{originalreference}{original reference}{original references}
\Crefname{originalreference}{Original reference}{Original references}
\makeatletter
\newlabel{more-algebra-lemma-quotient-sequence-regular}{{0GEE}{source}{}{sourcepointer-10}{}}
\newlabel{more-algebra-lemma-quotient-sequence-regular@cref}{{[originalreference][0][]0GEE}{[1][0][]source}{}{}{}}
\makeatother
\title{Rig-smooth implies rig-flat with regular sequence presentations}
\author{The Stacks Project Authors}
\date{stacks-project:restricted.tex:local-snapshot-20261012; GNU FDL 1.2 or later}
\begin{document}
\maketitle
\noindent\textit{Editorial note.} Counted span is the author's own proof environment, source lines 4470-4558, following the statement at 4459-4468.\par\medskip
\noindent\textit{Editorial note.} The parenthetical (details omitted) at line 4505 is the standard identity identifying restricted power series over the completed localization with the completed localization of restricted power series; recorded in the card's obligations.\par\medskip
\noindent\textit{Editorial note.} Steps the author leaves to the reader are flagged here for the reviewer: source line(s) 4505 carry an "omitted"/"follows easily"/"immediate from the definitions" phrase. Those sentences are retained verbatim; the displayed argument around them stays inside the counted span.\par\medskip
\noindent\textit{Allowed standard prerequisites (editorial).}\begin{itemize}
\item Rig-smooth base change and locality. lemma-base-change-rig-smooth-continuous and lemma-rig-smooth-axioms, cited at proof lines 4473-4477. Supplies that the localizations phi\_\{f\} stay rig-smooth, reducing rig-flatness to the naive condition.
\item Regular sequences on open subsets of V(J). algebra-lemma-regular-sequence-in-neighbourhood, cited at proof lines 4486-4489. Supplies that it suffices to check J\_q generated by a regular sequence at primes maximal for not containing an ideal of definition.
\item Rig-closed points after localization and descent. lemma-rig-closed-point-after-localization (4498-4502), algebra-lemma-flat-increases-depth and more-algebra-lemma-flat-descent-regular-ideal (4527-4528). Supplies the rig-closed prime p' and the faithfully flat descent of regular-sequence generation and flatness along the displayed diagram.
\item Fibre regularity and Grothendieck regular sequence criterion. lemma-fibre-regular (4545-4548), more-algebra-lemma-quotient-sequence-regular (4550-4552), algebra-lemma-grothendieck-regular-sequence (4553-4555). Supplies that f\_1,...,f\_m map to a regular sequence in the fibre, hence in the local ring, with flat quotient.
\end{itemize}\medskip
\ifdefined\hypertarget\hypertarget{sourcepointer-1}{}\fi
\noindent\textit{Source reference.} Original reference 00LM: focused-seeds/geometry/S9-restricted-rig-smooth-rig-flat/ref-tags.txt, source lines 5--5. Navigation pointer only: the counted proof cites this statement by \textbackslash{}ref as a standing input; the deduction the counted proof has to make is carried out inside the counted span, and no part of that deduction is delegated to this pointer.\par\medskip
\ifdefined\hypertarget\hypertarget{sourcepointer-2}{}\fi
\noindent\textit{Source reference.} Original reference 00MG: focused-seeds/geometry/S9-restricted-rig-smooth-rig-flat/ref-tags.txt, source lines 6--6. Navigation pointer only: the counted proof cites this statement by \textbackslash{}ref as a standing input; the deduction the counted proof has to make is carried out inside the counted span, and no part of that deduction is delegated to this pointer.\par\medskip
\ifdefined\hypertarget\hypertarget{sourcepointer-3}{}\fi
\noindent\textit{Source reference.} Original reference 061L: focused-seeds/geometry/S9-restricted-rig-smooth-rig-flat/ref-tags.txt, source lines 7--7. Navigation pointer only: the counted proof cites this statement by \textbackslash{}ref as a standing input; the deduction the counted proof has to make is carried out inside the counted span, and no part of that deduction is delegated to this pointer.\par\medskip
\ifdefined\hypertarget\hypertarget{sourcepointer-4}{}\fi
\noindent\textit{Source reference.} Original lemma 0GCL: sources/stacks-restricted/pinned-original-source.tex, source lines 4371--4371. Navigation pointer only: the counted proof cites this statement by \textbackslash{}ref as a standing input; the deduction the counted proof has to make is carried out inside the counted span, and no part of that deduction is delegated to this pointer.\par\medskip
\ifdefined\hypertarget\hypertarget{sourcepointer-5}{}\fi
\noindent\textit{Source reference.} Original lemma 0GAJ: sources/stacks-restricted/pinned-original-source.tex, source lines 843--843. Navigation pointer only: the counted proof cites this statement by \textbackslash{}ref as a standing input; the deduction the counted proof has to make is carried out inside the counted span, and no part of that deduction is delegated to this pointer.\par\medskip
\ifdefined\hypertarget\hypertarget{sourcepointer-6}{}\fi
\noindent\textit{Source reference.} Original lemma 0GGJ: sources/stacks-restricted/pinned-original-source.tex, source lines 3604--3604. Navigation pointer only: the counted proof cites this statement by \textbackslash{}ref as a standing input; the deduction the counted proof has to make is carried out inside the counted span, and no part of that deduction is delegated to this pointer.\par\medskip
\ifdefined\hypertarget\hypertarget{sourcepointer-7}{}\fi
\noindent\textit{Source reference.} Original lemma 0GGG: sources/stacks-restricted/pinned-original-source.tex, source lines 3488--3488. Navigation pointer only: the counted proof cites this statement by \textbackslash{}ref as a standing input; the deduction the counted proof has to make is carried out inside the counted span, and no part of that deduction is delegated to this pointer.\par\medskip
\ifdefined\hypertarget\hypertarget{sourcepointer-8}{}\fi
\noindent\textit{Source reference.} Original lemma 0GCK: sources/stacks-restricted/pinned-original-source.tex, source lines 4253--4253. Navigation pointer only: the counted proof cites this statement by \textbackslash{}ref as a standing input; the deduction the counted proof has to make is carried out inside the counted span, and no part of that deduction is delegated to this pointer.\par\medskip
\ifdefined\hypertarget\hypertarget{sourcepointer-9}{}\fi
\noindent\textit{Source reference.} Original reference 068N: focused-seeds/geometry/S9-restricted-rig-smooth-rig-flat/ref-tags.txt, source lines 8--8. Navigation pointer only: the counted proof cites this statement by \textbackslash{}ref as a standing input; the deduction the counted proof has to make is carried out inside the counted span, and no part of that deduction is delegated to this pointer.\par\medskip
\ifdefined\hypertarget\hypertarget{sourcepointer-10}{}\fi
\noindent\textit{Source reference.} Original reference 0GEE: focused-seeds/geometry/S9-restricted-rig-smooth-rig-flat/ref-tags.txt, source lines 9--9. Navigation pointer only: the counted proof cites this statement by \textbackslash{}ref as a standing input; the deduction the counted proof has to make is carried out inside the counted span, and no part of that deduction is delegated to this pointer.\par\medskip
\begin{lemma}
\label{lemma-rig-smooth-rig-flat}
Let $\varphi : A \to B$ be an arrow of $\textit{WAdm}^{Noeth}$.
If $\varphi$ is rig-smooth, then $\varphi$ is rig-flat, and
for any presentation $B = A\{x_1, \ldots, x_n\}/J$
and prime $J \subset \mathfrak q \subset A\{x_1, \ldots, x_n\}$
not containing an ideal of definition the ideal
$J_\mathfrak q \subset A\{x_1, \ldots, x_n\}_\mathfrak q$
is generated by a regular sequence.
\end{lemma}

\begin{proof}
Let $f \in A$. To prove that $\varphi$ is rig-flat we have to show
that $\varphi_{\{f\}} : A_{\{f\}} \to B_{\{f\}}$ is naively rig-flat.
Now either by viewing $\varphi_{\{f\}}$ as a base change of $\varphi$
and using Lemma \ref{lemma-base-change-rig-smooth-continuous}
or by using the fact that being rig-smooth
is a local property (Lemma \ref{lemma-rig-smooth-axioms}) we see that
$\varphi_{\{f\}}$ is rig-smooth. Hence it suffices to show
that $\varphi$ is naively rig-flat.

\medskip\noindent
Choose a presentation $B = A\{x_1, \ldots, x_n\}/J$.
In order to check the second part of the lemma it suffices
to check $J_\mathfrak q \subset A\{x_1, \ldots, x_n\}_\mathfrak q$
is generated by a regular sequence for $J \subset \mathfrak q$
for $\mathfrak q$ maximal with respect to not containing
an ideal of definition, see
Algebra, Lemma \ref{algebra-lemma-regular-sequence-in-neighbourhood}
(which shows that the set of primes in $V(J)$ where there is
a regular sequence generating $J$ is open).
In other words, we may assume $\mathfrak q$ is rig-closed
in $A\{x_1, \ldots, x_n\}$. And to check that
$B$ is naively rig-flat, it also suffices to
check that the corresponding localizations $B_\mathfrak q$
are flat over $A$.

\medskip\noindent
Let $\mathfrak q \subset A\{x_1, \ldots, x_n\}$ be rig-closed with
$J \subset \mathfrak q$. By
Lemma \ref{lemma-rig-closed-point-after-localization}
we may choose an $f \in A$ mapping to a unit in
$A\{x_1, \ldots, x_n\}/\mathfrak q$ and such
that the prime ideal $\mathfrak p'$ in $A_{\{f\}}$ induced is rig-closed.
Below we will use that
$A_{\{f\}}\{x_1, \ldots, x_n\} = A\{x_1, \ldots, x_n\}_{\{f\}}$;
details omitted. Consider the diagram
$$
\xymatrix{
A\{x_1, \ldots, x_n\}_{\mathfrak q} / J_\mathfrak q \ar[r] &
A_{\{f\}}\{x_1, \ldots, x_n\}_{\mathfrak q'}/
J A_{\{f\}}\{x_1, \ldots, x_n\}_{\mathfrak q'} \\
A\{x_1, \ldots, x_n\}_{\mathfrak q} \ar[r] \ar[u] &
A_{\{f\}}\{x_1, \ldots, x_n\}_{\mathfrak q'} \ar[u] \\
A \ar[r] \ar[u] &
A_{\{f\}} \ar[u]
}
$$
The middle horizontal arrow is faithfully flat as it is a local
homomorphism of local rings and flat as $A_{\{f\}}\{x_1, \ldots, x_n\}$
is the completion of a localization of the Noetherian ring
$A\{x_1, \ldots, x_n\}$. Similarly the bottom horizontal arrow is flat.
Hence to show that $J_\mathfrak q$ is generated by a regular sequence
and that $A \to A\{x_1, \ldots, x_n\}_{\mathfrak q} / J_\mathfrak q$
is flat, it suffices to prove the same things for
$J A_{\{f\}}\{x_1, \ldots, x_n\}_{\mathfrak q'}$ and
$A_{\{f\}} \to A_{\{f\}}\{x_1, \ldots, x_n\}_{\mathfrak q'}/
J A_{\{f\}}\{x_1, \ldots, x_n\}_{\mathfrak q'}$.
See Algebra, Lemma \ref{algebra-lemma-flat-increases-depth} or
More on Algebra, Lemma \ref{more-algebra-lemma-flat-descent-regular-ideal}
for the statement on regular sequences. Finally, we have already seen that
$A_{\{f\}} \to B_{\{f\}}$ is rig-smooth.
This reduces us to the case discussed in the next paragraph.

\medskip\noindent
Let $\mathfrak q \subset A\{x_1, \ldots, x_n\}$ be rig-closed with
$J \subset \mathfrak q$ such that moreover $\mathfrak p = A \cap \mathfrak q$
is rig-closed as well. By the characterization of rig-smooth algebras
given in Lemma \ref{lemma-equivalent-with-artin-smooth}
after reordering the variables $x_1, \ldots, x_n$
we can find $m \geq 0$ and $f_1, \ldots, f_m \in J$ such that
\begin{enumerate}
\item $J_\mathfrak q$ is generated by $f_1, \ldots, f_m$, and
\item $\det_{1 \leq i, j \leq m}(\partial f_j/ \partial x_i)$
maps to a unit in $A\{x_1, \ldots, x_n\}_\mathfrak q$.
\end{enumerate}
By Lemma \ref{lemma-fibre-regular} the fibre ring
$$
F = A\{x_1, \ldots, x_n\} \otimes_A \kappa(\mathfrak p)
$$
is regular. Observe that the $A$-derivations $\partial / \partial x_i$
extend (uniquely) to derivations $D_i : F \to F$. By
More on Algebra, Lemma \ref{more-algebra-lemma-quotient-sequence-regular}
we see that $f_1, \ldots, f_m$ map to a regular sequence in
$F_\mathfrak q$. By flatness of $A \to A\{x_1, \ldots, x_n\}$
and Algebra, Lemma \ref{algebra-lemma-grothendieck-regular-sequence}
this shows that $f_1, \ldots, f_m$ map to a regular sequence in
$A\{x_1, \ldots, x_m\}_\mathfrak q$ and the quotient by
these elements is flat over $A$. This finishes the proof.
\end{proof}

\end{document}
